\chapter{Build, test and delivery}\label{ch:delivery}
\chapterstatus{complete}{2026-11-16}{\src{.github/workflows/deploy.yml}, \src{infra/aws/deploy.sh}, \src{Dockerfile}}

Every push and pull request is checked; every push to \code{master} is deployed. Nobody deploys
by hand and nobody holds long-lived AWS keys for it.

\begin{figure}[htbp]
  \centering
  % CI/CD pipeline (chapter 9): .github/workflows/ci.yml, .github/workflows/deploy.yml, infra/aws/deploy.sh.
\begin{tikzpicture}[node distance=4mm and 4mm,
    st/.style={svc=23mm, minimum height=11mm, font=\sffamily\scriptsize},
    gate/.style={st, draw=ink, fill=wash}]
  % CI: every pull request
  \node[ext, text width=16mm] (pr) {pull request};
  \node[st, text width=50mm, right=of pr] (jobs) {CI jobs, in parallel\\[-1pt]\textcolor{muted}{\tiny static + unit · integration · end-to-end · images · docs}};
  \node[gate, right=of jobs] (ok) {CI ok\\[-1pt]\textcolor{muted}{\tiny required, up to date}};
  \draw[flow] (pr) -- (jobs);
  \draw[flow] (jobs) -- (ok);
  % Deploy: every merge
  \node[ext, text width=16mm, below=9mm of pr] (merge) {merge to master};
  \node[st, right=of merge] (build) {build ARM64\\[-1pt]\textcolor{muted}{\tiny app + migrate, one bake}};
  \node[aws=compute, box=20mm, minimum height=11mm, right=of build] (ecr) {ECR\\[-1pt]\textcolor{muted}{\tiny <sha>}};
  \node[gate, right=of ecr] (mig) {migrate task\\[-1pt]\textcolor{muted}{\tiny if migrations changed}};
  \node[st, below=6mm of mig] (roll) {update services\\[-1pt]\textcolor{muted}{\tiny circuit breaker}};
  \node[gate, left=of roll] (verify) {verify rollout\\[-1pt]\textcolor{muted}{\tiny new release on /health}};
  \draw[flow] (ok.south) -- ++(0,-3mm) -| node[edgelabel, pos=0.25, above]{merge} (merge.north);
  \draw[flow] (merge) -- (build);
  \draw[flow] (build) -- (ecr);
  \draw[flow] (ecr) -- (mig);
  \draw[flow] (mig) -- (roll);
  \draw[flow] (roll) -- (verify);
\end{tikzpicture}

  \caption{The pipeline. A failing step stops everything after it; production keeps the previous release.}
  \label{fig:pipeline}
\end{figure}

\section{Checks}
The \code{check} job runs on every push and pull request: \code{bun install --frozen-lockfile},
\code{bun run typecheck}, \code{bun run lint}, \code{bun run test:unit}. Typecheck and lint cover the
project's own code and skip the AI-DLC folders \src{.claude} and \src{aidlc} (\cref{ch:aidlc}), which are
vendored and replaced on each framework update. Integration and end-to-end
tests need PostgreSQL and run locally (\code{bun run test}, \code{bun run test:e2e}); see
\cref{sec:local-checks}.

\section{Images}
The \code{deploy} job builds two images natively on GitHub's ARM64 runner, matching Fargate's
Graviton CPUs, and pushes them to ECR tagged with the commit SHA:

\begin{xltabular}{\linewidth}{@{}P{30mm} P{28mm} L@{}}
  \toprule
  \thead{Image} & \thead{Tag} & \thead{Contents}\\
  \midrule
  App (default target) & \code{<sha>} & Next.js standalone server, the worker and admin scripts bundled by Bun, the RDS CA bundle. About 450 MB.\\
  Migrate (\code{--target migrate}) & \code{<sha>-migrate} & The app image plus the Prisma CLI and OpenSSL; trusts only the RDS CA bundle.\\
  \bottomrule
\end{xltabular}

Tags are immutable, so a re-run of the same commit reuses its images instead of rebuilding.

\section{Rollout}
\src{infra/aws/deploy.sh} does the rest, and works the same by hand:
\begin{enumerate}
  \item Registers new revisions of the \code{slotty-migrate}, \code{slotty-web} and
    \code{slotty-worker} task definitions: the latest revision with only the image changed.
    Configuration changes made in Terraform therefore ride along with the next deploy.
  \item Runs the migrate task, waits for it to stop and prints its log. A non-zero exit stops the
    deploy; the services stay on the previous release.
  \item Updates the web and worker services and waits until they are stable.
  \item Waits for each new deployment's own verdict (completed or failed, up to 10 minutes) and
    checks it is the one now running. A circuit-breaker rollback is reported as a failure.
  \item Requests \route{/health} through CloudFront, which also checks the database.
\end{enumerate}

\begin{decision}{ADR-010}
  CI owns which image runs; Terraform owns everything else. ECS services ignore task-definition
  changes in Terraform, and CI's role can only push images, register task definitions, run the
  migrate task and update the two services.
\end{decision}

\section{Identity}
GitHub Actions assumes the \code{slotty-github-deploy} role with OpenID Connect. The role trusts
only tokens whose subject is this repository's \code{master} branch, in GitHub's immutable form
(\code{repo:<owner>@<owner-id>/<repo>@<repo-id>:ref:refs/heads/master}), so a renamed or
re-created repository with the same name cannot deploy. Its permissions are listed in
\src{infra/aws/app/iam.tf}.

The workflow reads its targets from repository variables (\code{AWS_DEPLOY_ROLE_ARN},
\code{ECR_REPOSITORY}, \code{ECS_CLUSTER}, \code{ECS_SUBNETS}, \code{ECS_SECURITY_GROUP},
\code{APP_URL}), which \code{terraform output github_actions_variables} prints. None of them are
secrets. Until \code{AWS_DEPLOY_ROLE_ARN} is set, only the checks run.

\section{Local checks}\label{sec:local-checks}
\begin{lstlisting}[language=bash]
bun run infra:up        # Postgres, Mailpit, Caddy (Docker)
bun run typecheck && bun run lint && bun run test && bun run test:e2e
bun run app:up          # the production images locally at http://localhost:3001
\end{lstlisting}

\section{This document}
\code{bun run docs:build} builds this PDF in a pinned TeX Live image (\src{docs/tooling}). CI
builds it for every pull request and keeps the PDF as a build artifact, so a broken document
fails like broken code.
